\section{Methods}

\subsection{Cohort provenance, skip-logic reconstruction, and variable audit}
We analyzed all 368,132 isolated-CABG procedures in STS ACSD v2.81 (July 1, 2014--June 30, 2017; 1,091 US surgical centers across all nine US Census divisions)\cite{sts2014v281,shahian2018sts} under institutional Data Use Agreement governance (IRB-exempt de-identified secondary research), following TRIPOD+AI (Supplementary Table~S3) and PROBAST+AI guidelines.\cite{collins2024tripodai,moons2025probastai}

Three schema audits preceded modeling (Supplementary Section~S1). First, skip-logic outcome reconstruction (Sequence 4930) resolved structurally blank \cli{COtAFib} entries among 232,783 uncomplicated recoveries (\cli{Complics == 2}, 63.31\%) as POAF-negative to prevent collider bias,\cite{hernan2020causal} recovering \textbf{367,695 labeled cases} (437 genuinely unknown excluded; \textbf{89,337 POAF events, 24.30\%}). Second, column-by-column auditing purged 431 post-incision or identifier variables---including seven early-column postoperative duration leaks (\cli{LOS}, \cli{PLOS}, \cli{ORDuration}, \cli{SIDuration}, \cli{SurgtoMort}, \cli{SurgtoRead}, \cli{SurgtoSternInf}), cardiopulmonary bypass/cross-clamp durations, discharge medications, STS risk calculators, auxiliary outcomes, and hospital/surgeon identifiers---retaining strictly \textbf{207 raw pre-incision predictors} (including pre-incision operating-room \cli{PASYS} where measured prior to skin incision and 9-level \cli{Region}). Third, sentinel non-numeric codings (e.g., 999) were mapped to missing values prior to model fitting.

\subsection{Dual unseen-hospital and seen-hospital split design}
To separate cold-start external generalization from in-network hospital deployment, we constructed two matched evaluation splits with identical natural-prior holdout size ($N{=}21{,}241$, POAF prevalence $\approx$24.3\%) and identical 50,000-patient balanced five-fold hospital-grouped training cohorts (Table~\ref{tab:cohort}):
\begin{enumerate}
  \item \textbf{Regime 1 (Unseen Hospitals, 0.0\% overlap):} Multivariate anticlustering partitioned the 1,091 hospitals into 1,017 development centers (346,454 cases; nested 50,000-patient five-fold hospital-disjoint benchmark) and \textbf{74 strictly unseen holdout hospitals} ($N{=}21{,}241$, 5,160 events, $\mathrm{Max}\,|\mathrm{SMD}|=0.103$).\cite{papenberg2021anticlustering,joseph2022split}
  \item \textbf{Regime 2 (Seen Hospitals, 100.0\% overlap):} Hospital-stratified splitting ensured that \textbf{all 1,091 surgical centers} appeared in both the 346,454-case development cohort and the nested 50,000-patient training cohort, with a matched holdout ($N{=}21{,}241$ across 1,086 seen centers, zero unseen centers). Models were scored first using only the 207 patient-level predictors (zero site ID) and second with strictly train-only Empirical-Bayes (EB) hospital POAF priors via both closed-form log-odds site calibration ($\mathrm{EB}_{50\mathrm{k}}$, $m_0{=}15$; $\mathrm{EB}_{346\mathrm{k}}$, $m_0{=}25$) and native end-to-end 208-feature training (\cli{site208}).
\end{enumerate}

\subsection{Binary comparator roster, proprietary benchmark, and latency protocol}
Five leading public tabular AI systems were trained strictly on binary $\{0,1\}$ POAF labels across identical 50,000-patient five-fold hospital-grouped cohorts under active Python audit hooks blocking test-label access: (1) \textbf{Google TabFM v1.0.0}\cite{kong2026tabfm} ($5\times 10\text{k}$ bagged ensemble, 80 estimators); (2) \textbf{NVIDIA KumoTabular-Large v2}\cite{qu2026kumotabular} ($5\times 10\text{k}$ bagged ensemble, 40 estimators with train-only class-orientation checks); (3) \textbf{PriorLabs TabPFN-3.5}\cite{hollmann2025tabpfn,jaeger2026tabpfn35} (native full 50,000-row in-context classification, 8 estimators); (4) \textbf{CatBoost}\cite{prokhorenkova2018catboost} (tuned five-fold symmetric trees, depth 6, plus a $7\times$ 346,454-patient full-registry scaling reference); and (5) \textbf{AutoGluon-Tabular 1.5}\cite{erickson2020autogluon} (\cli{best\_quality} multi-layer stack of LightGBM, CatBoost, XGBoost, PyTorch/FastAI neural networks, random forests, and extra trees, strictly enforcing our five hospital-disjoint folds via \cli{groups="cv\_fold"}, 30-minute training budget per split on GPU+CPU). These five alternatives were benchmarked head-to-head against \textbf{NEVERHUMAN}, a proprietary pre-incision clinical risk model restricted to the exact same 207 pre-incision inputs, training cohorts, and evaluation pipeline, alongside a pre-incision reconstruction of the bedside Mariscalco POAF score.\cite{mariscalco2014poaf}

All 50,000-patient models received uniform fold-out-of-fold Platt scaling\cite{platt1999probabilistic} plus development-only Bayes prior shift ($-1.1365$) to restore natural prevalence.\cite{king2001logistic,saerens2002adjusting} Training/preparation wall-clock time, full-cohort inference wall-clock time ($N{=}21{,}241$; median of three passes for deployed models and verified offline GPU receipts for foundation models), per-patient latency ($\mathrm{ms/patient}$), and peak GPU memory were benchmarked cleanly across architectures.

\subsection{Statistical analysis and clinical utility}
All confidence intervals and paired model comparisons used exact \textbf{5,000-draw whole-hospital cluster bootstrapping} (seed 42; 74 unseen centers / 1,086 seen centers).\cite{davison1997bootstrap,obuchowski1997clustered} Beyond ROC-AUC and average precision, we evaluated 15-bin expected calibration error (ECE), calibration slope, observed-to-expected (O:E) ratio,\cite{vancalster2019calibration,austin2019ici} Decision Curve Analysis (DCA) net benefit and avoided unnecessary interventions per 100 patients across 15\%--40\% thresholds,\cite{vickers2006decision} and 25 prespecified demographic, clinical, acuity, and hospital-volume subgroups.
